\subsection{APPLICATION : THE POINCARÉ-VOLTERRA THEOREM}

THEOREM 1. Let X be a topological space satisfying axiom (O \(_{III}\) ) (but not necessarily Hausdorff), and suppose X is connected and locally connected. Let Y be a topological space whose topology has a countable base, and let p: X → Y be a continuous mapping such that, for each y ∈ Y, \(\overline{p}^{1}(y)\) is a discrete subspace of X. Finally let B be a set of subsets of X whose interiors cover X and such that:

\begin{enumerate}
\def\labelenumi{(\roman{enumi})}
\tightlist
\item
  The restriction of \(p\) to each \(V \in \mathfrak{B}\) is a closed mapping of \(V\) into \(Y\) .\\
\item
  Every \(V \in \mathfrak{B}\) has a countable subset which is dense in \(V\) .
\end{enumerate}

Then the space X is the union of a countable family of open sets each of which is contained in a set of B.

Let \(\mathfrak{B}\) be a countable base of the topology of Y. We shall say that a pair (W, U) is distinguished if (i) \(U \in \mathfrak{B}\) and (ii) W is a component of \(\overline{p}(U)\) contained in a set of \(\mathfrak{B}\) .

LEMMA 1. If \(x\) is any point of \(\mathbf{X}\) , there is a distinguished pair \((\mathbf{W}, \mathbf{U})\) such that \(x \in \mathbf{W}\) .

The inverse image \(\overline{p} (p(x))\) is discrete and therefore there is a neighbourhood of \(x\) in \(\mathbf{X}\) all of whose points \(x^{\prime}\) other than \(x\) have an image \(p(x^{\prime})\neq p(x)\) ; since \(\mathbf{X}\) satisfies \((\mathrm{O}_{\mathrm{III}})\) , there is a closed neighbourhood \(\mathbf{V}\) of \(x\) with this property, and we may assume also that \(\mathbf{V}\) is contained in a set of \(\mathfrak{B}\) . Let \(\mathbf{F}\) be the frontier of \(\mathbf{V}\) in \(\mathbf{X}\) . By condition (i) of the theorem, \(p(\mathbf{F})\) is closed in \(\mathbf{Y}\) ; and since \(p(\mathbf{F})\) does not contain \(p(x)\) , there is a set \(\mathbf{U}\in \mathfrak{B}\) which contains \(p(x)\) and does not meet \(p(\mathbf{F})\) . Let \(\mathbf{W}\) be the component of \(x\) in \(\overline{p}^{1}(\mathbf{U})\) ; then it is enough to show that \(\mathbf{W}\subset \mathring{\mathbf{V}}\) . If this were not so, then \(\mathbf{W}\) would meet \(\mathbf{F}\) (no. 1, Proposition 3) and therefore \(p(\mathbf{F})\) would meet \(\mathbf{U}\) , contrary to the definition of \(\mathbf{U}\) .

LEMMA 2. If (W, U) is a distinguished pair then the set of all distinguished pairs (W', U') such that W' meets W is countable.

Since \(\mathfrak{B}\) is countable it is enough to show that, given \(\mathbf{U}' \in \mathfrak{B}\) , the set of distinguished pairs \((\mathbf{W}', \mathbf{U}')\) such that \(\mathbf{W}'\) meets \(\mathbf{W}\) is countable. Now these sets \(\mathbf{W}'\) are open, since \(\mathbf{X}\) is locally connected (no. 6, Proposition 11) and mutually disjoint since they are components of \(\overline{\boldsymbol{p}}(\mathbf{U}')\) ; hence the sets \(\mathbf{W}' \cap \mathbf{W}\) are open and mutually disjoint. But \(\mathbf{W}\) contains a countable subset which is dense in \(\mathbf{W}\) ; hence the set of \(\mathbf{W}'\) such that \(\mathbf{W}' \cap \mathbf{W}\) is not empty is also countable.

To prove Theorem 1, consider the following relation R between two points \(x\) , \(x'\) of X: ``There exists a finite sequence of distinguished pairs \((W_i, U_i)\) ( \(1 \leqslant i \leqslant n\) ) such that \(x \in W_1\) and \(x' \in W_n\) and that \(W_i \cap W_{i+1} \neq \emptyset\) for \(1 \leqslant i \leqslant n - 1\) .''

Lemma 1 states that R is reflexive, and it is readily verified that R is symmetric and transitive, so that R is an equivalence relation; also, since the \(W_{i}\) are open, every equivalence class mod R is open in X. But X is connected; hence there can be only one equivalence class, i.e.~any two points of X are congruent mod R. We shall deduce from this that X is the union of a countable family of first elements of distinguished pairs, and this will prove Theorem 1. For this, let x be any point of X, and define by induction on n a sequence \((\mathbf{C}_{n})\) of open subsets of X as follows: by Lemma 1 there is a distinguished pair \((\mathbf{W}_{1}, \mathbf{U}_{1})\) such that \(x \in W_{1}\) , and we take \(C_{1} = W_{1}\) ; if n \textgreater{} 1 then \(C_{n}\) is to be the union of all first elements W of distinguished pairs (W, U) such that W meets \(C_{n-1}\) . By induction on n one shows immediately, by virtue of Lemma 2, that \(C_{n}\) is a countable union of first elements of distinguished pairs. Finally, every \(x' \in X\) belongs to some \(C_{n}\) ; for there is a finite sequence

\[
\left(\mathrm{W} _ {i} ^ {\prime}, \mathrm{U} _ {i} ^ {\prime}\right) _ {1 \leqslant i \leqslant m}
\]

of distinguished pairs such that \(x \in W_1'\) , \(x' \in W_m'\) and \(W_i' \cap W_{i+1}' \neq \emptyset\) for \(1 \leqslant i \leqslant m - 1\) , and by induction on \(i\) we see that \(W_i' \subset C_{i+1}\) for all \(i\) , so that \(x' \in C_{m+1}\) .

Q.E.D.

COROLLARY 1. Let Y be a regular space whose topology has a countable base (*). Let X be a connected and locally connected space, and let \(p: \mathbf{X} \to \mathbf{Y}\) be a continuous mapping with the following property: for each \(x \in \mathbf{X}\) there is a closed neighbourhood V of \(x\) in X such that the restriction of \(p\) to V is a homeomorphism of V onto a closed subspace of Y. Then X is regular and the topology of X has a countable base.

First, the hypotheses imply that X is regular (§ 8, no. 4, Proposition 13). Let us show that the conditions of Theorem 1 are satisfied if we take B to be the set of all closed subsets V of X such that the restriction of p to V is a homeomorphism of V onto a closed subspace of Y. By hypothesis, the interiors of the sets of B cover Y and, by virtue of the assumption on Y, each \(V \in B\) has a countable base and therefore contains a countable dense subset (§ 1, no. 6, Proposition 6). Furthermore, if \(x \in \overline{p}(y)\) , there is a neighbourhood \(V \in B\) of x in X such that V contains no point of \(\overline{p}(y)\) other than \(x\) , and hence \(\overline{p}(y)\) is a discrete space. We may therefore apply Theorem 1, which shows that X is the union of a countable family \((\mathrm{T}_n)_{n\geqslant 0}\) of open sets, such that each subspace \(\mathrm{T}_n\) has a countable base \((\mathrm{U}_{mn})_{m\geqslant 0}\) . Then the \(\mathrm{U}_{mn}\) ( \(m\geqslant 0\) , \(n\geqslant 0\) ) form a base of the topology of X (§ 3, no. 1, Remark).

COROLLARY 2. Let X be locally compact, connected and locally connected, and suppose each point of X has a neighbourhood which has a countable base. Let Y be a Hausdorff space whose topology has a countable base, and let \(p: \mathbf{X} \to \mathbf{Y}\) be a continuous mapping such that, for each \(y \in \mathbf{Y}\) , \(\overline{p}^1(y)\) is a discrete subspace of X. Then the topology of X has a countable base.

For each \(x \in X\) , let \(V_{x}\) be a compact neighbourhood of x in X which has a countable base. It follows from §9, no. 4, Theorem 2, Corollary 2, that the set B of the \(V_{x}\) satisfies the conditions of Theorem 1, and we complete the proof as in Corollary 1.

Notice that, in this corollary, it can happen that the restriction of p to an arbitrarily small neighbourhood V of a point of X is not a homeomorphism of V onto \(p(\mathrm{V})\) .

COROLLARY 3 (Poincaré-Volterra Theorem). Let Y be a locally compact, locally connected space whose topology has a countable base. Let X be a connected Hausdorff space, and let \(p: \mathbf{X} \to \mathbf{Y}\) be a continuous mapping which has the following property: for each \(x \in \mathbf{X}\) there is an open neighbourhood U of \(x\) in X such that the restriction of \(p\) to U is a homeomorphism of U onto an open subspace of Y. Then X is locally compact and locally connected, and the topology of X has a countable base.

Clearly X is locally connected. Also each \(x \in X\) has an open neighbourhood U in X such that the restriction of p to U maps U homeomorphically onto an open subspace \(p(U)\) of Y. Since \(p(U)\) is a locally compact subspace of Y ( \(§ 9, no. 7, Proposition 13\) ), there is a compact neighbourhood W of \(p(x)\) contained in \(p(U)\) , whence \(U \cap \overline{p}(W)\) is a compact neighbourhood of x contained in U; thus X is locally compact, since by hypothesis X is Hausdorff. \(U \cap \overline{p}(W)\) , being compact, is closed in X ( \(§ 9, no. 3, Proposition 4\) ) and the conditions of Corollary 1 are therefore satisfied; hence the topology of X has a countable base.

